\section{A probability flow in the subexponential case}
\label{sec:flow}

Assume now that the unbounded moduli satisfy
\begin{equation}
 \limsup_{j\to\infty}M(j)^{1/j}\leq1.
 \label{eq:subexponential-growth}
\end{equation}
Let $u,\mu,C_t$ be given by \Cref{prop:scalar-limit}. We build a tree of
small balls carrying nested intervals of values of $u$. At each node,
either two separated variations survive expansion, or many children
carry a substantial fraction of its interval. A probability flow will
force the first alternative at arbitrarily deep levels.

Choose $x_0\in U$ where $u$ is not locally constant, and choose
$\epsilon_0>0$ with $\overline B(x_0,3\epsilon_0)\subset U$. Put
\[
 K=\overline B(x_0,\epsilon_0),\qquad
 K^+=\overline B(x_0,2\epsilon_0),
\]
and, shrinking $\epsilon_0$ if necessary, choose $z_0$ a positive distance
from $K^+$. All centers will lie in $K$ and all source neighborhoods in
$K^+$. Use the constants of \eqref{eq:measure} for this fixed region.

\subsection{Node intervals and transition probabilities}

The tree will record both where variation occurs and which values it
carries. A node has a source center $x$, a scale $R$, and a value interval
$J$. Its children have scale $r=\theta R$, where $\theta=a^{-D}$ for
one positive integer subdivision depth $D$ used throughout the tree. Their intervals
are contained in $J$. We will choose probabilities on the children
that control interval length linearly and quadratically; the two
controls have different roles in the final comparison.

We first define this data for geometric constants $c_0>2L+2$,
$Lc_0+2<A<B_0-2$, and $H>0$. Here $c_0$ is the source-ball enlargement,
$A,B_0$ locate an annulus of possible child centers, and $Hr$ is the
separation required of retained children. The next subsection fixes
these constants, before choosing $D$. We use only scales $R=r_k$
small enough that all fixed enlargements below lie in $K^+$ and
satisfy \eqref{eq:measure}, and that $(B_0+2)R<d_1/2$.

A \emph{node} is a center $x\in K$, a scale $R=r_k$, and a nonempty
open interval $J=(j_-,j_+)\subset(0,1)$ such that
\begin{equation}
 J\subset\left(\min_{\overline B(x,c_0R)}u,
                   \max_{\overline B(x,c_0R)}u\right).
 \label{eq:node-straddling}
\end{equation}
Write $\Delta=|J|$. For every $t\in J$, a path in $B(x,Lc_0R)$
strictly straddles $t$ and hence meets $C_t$. Since
$\diam C_t\geq d_1$, the continuum leaves
$\overline B(x,(B_0+2)R)$ and meets every intervening radial sphere.

Take a maximal $r$-separated set $Y$ in
\[
 \mathcal A_x(R)=\{y:(A-1)R\leq d(x,y)\leq(B_0+1)R\}.
\]
It is an $r$-net there, and packing gives
\begin{equation}
 \#Y\leq C_Ya^{QD},
 \label{eq:candidate-count}
\end{equation}
with $C_Y$ independent of $D$. For $y\in Y$ let $I_y^{\mathrm{in}}$
and $I_y^{\mathrm{out}}$ be the closed intervals between the extrema
of $u$ on $\overline B(y,2r)$ and $\overline B(y,c_0r)$, respectively.
These are range \emph{intervals}; the balls need not be connected. Put
\[
 J_y=J\cap\operatorname{int}I_y^{\mathrm{out}},\qquad
 t_y=|J_y|/\Delta.
\]
A positive $t_y$ gives a nested open child interval satisfying
\eqref{eq:node-straddling}, and, including singleton ranges,
\begin{equation}
 |J\cap I_y^{\mathrm{out}}|=t_y\Delta.
 \label{eq:closed-range-length}
\end{equation}

For $0<c_*\leq1/(8\max\{C_Y,1\})$, put $\kappa=c_*a^{-QD}$.
A node is \emph{good} if two candidates $y^-,y^+$ satisfy
\begin{align}
 |I_{y^-}^{\mathrm{in}}\cap(j_-,j_-+\Delta/4)|&\geq\kappa\Delta,
 &|I_{y^+}^{\mathrm{in}}\cap(j_-+3\Delta/4,j_+)|&\geq\kappa\Delta,
 \label{eq:good-inner}\\
 \dist(I_{y^-}^{\mathrm{out}},I_{y^+}^{\mathrm{out}})&\geq\Delta/4.
 \label{eq:good-outer}
\end{align}
Goodness uses the full candidate set, independently of the children
later retained.

The required probabilities are as follows. Fix the exponents
\begin{equation}
 b_* =\frac12,\qquad \tau=\frac1{20},\qquad
 \beta=\frac{\tau}{4(Q+1)},\qquad \gamma=\frac\beta2.
 \label{eq:flow-exponents}
\end{equation}
At every node we seek an $Hr$-separated set of children, with
probabilities summing to one, such that
\begin{align}
 p_y&\leq a^{-\gamma D}t_y,
 &p_y&\leq a^{-b_*D}t_y^2
 &&\text{at nodes that are not good},\label{eq:transition-bad}\\
 p_y&\leq a^{-\gamma D}t_y,
 &p_y&\leq a^{\tau D}t_y^2
 &&\text{at good nodes}.\label{eq:transition-general}
\end{align}
Multiplying the linear bounds along a branch will telescope the
interval ratios and make the branch probability small. Multiplying
the quadratic bounds will compare probability to squared interval
length. The measure estimate bounds the sum of those squares at
each separated generation. It will therefore bound the average
probability of nodes with the stronger quadratic bound and force a
substantial amount of good-node probability at arbitrarily deep levels.

Two facts will supply these transitions. A node that is not good has
many children with $t_y\ge1/8$, on which uniform probabilities suffice.
At every node, a capacity estimate supplies enough total weight to
satisfy the weaker pair of bounds \eqref{eq:transition-general}.
After constructing the transitions, we show that goodness also gives
two separated regions of variation in one of finitely many target
patterns. That geometric property is the input to \Cref{sec:pairs}.

\subsection{Geometric choices and many large children}

We now fix the geometry used by the definitions above. Put
$c_{\mathrm{in}}=2L$ and choose $c_0>c_{\mathrm{in}}+2$ so large that,
with
\[
 A=4Lc_0+20,\qquad B_0=A+10,\qquad
 c_{\mathrm{pat}}=\frac{2C_Ec_{\mathrm{in}}}{(A-2)^2},
\]
we have
\begin{equation}
 \frac{c_Ec_0}{(B_0+2)^2}>12c_{\mathrm{pat}}.
 \label{eq:flow-constant-choice}
\end{equation}
The constant $c_{\mathrm{pat}}$ will be the target radius factor:
paths joining extrema of the two inner child balls will expand into balls of radius
$c_{\mathrm{pat}}\theta/2$, while the larger source balls will contain
the preimages of radius-$10c_{\mathrm{pat}}\theta$ target balls.
The ratio of the left side of \eqref{eq:flow-constant-choice} to the
right side grows linearly in $c_0$, so this choice is possible.
Next choose
\[
 H>\max\{4(B_0+1),4C_{\mu}c_0,2\}.
\]
The first inequality preserves separation between different branches;
the second makes the measure enlargements at one generation disjoint.
These choices fix $C_Y$ in \eqref{eq:candidate-count}, and we fix
$c_*$ as specified above. Require $\theta<1/4$ and $c_0\theta<1$.
All further lower bounds on $D$ will depend only on this fixed data.

The uniform upper bound on node scales can be chosen now, before $D$:
all source radius factors are fixed, and the subsidiary scales $r$
and $s$ used below never exceed $R$. Thus the source neighborhoods
and measure enlargements stay in $K^+$ for every node. We choose the
root scale after $D$ has been fixed.

\begin{lemma}[Many children at a node that is not good]
\label{lem:bad-node-children}
For sufficiently large $D$, every node that is not good has at least
$c_Ha^D$ mutually $Hr$-separated candidates with $t_y\geq1/8$,
where $c_H>0$ is independent of the node and $D$.
\end{lemma}
\begin{proof}
Fix $b\in[AR,B_0R]$. Each $C_t$, $t\in J$, meets the radius-$b$
sphere, so the inner intervals of candidates with
$|d(x,y)-b|\leq r$ cover $J$. By \eqref{eq:candidate-count}, two of
them meet the respective quarters in \eqref{eq:good-inner} by at least
$\kappa\Delta$.

If both have $t_y<1/8$, the lower-quarter outer interval ends below
$j_-+3\Delta/8$, and the upper-quarter interval starts above
$j_-+5\Delta/8$. Indeed each meets its quarter but overlaps $J$ in
length less than $\Delta/8$; the assertion also holds for intervals
extending beyond $J$. Their distance exceeds $\Delta/4$, making the
node good. Otherwise one candidate has $t_y\geq1/8$. Apply this at
radial values spaced by $(H+3)r$: there are at least $c_Ha^D$ such
values, and the resulting centers are more than $Hr$ apart.
\end{proof}

At a node that is not good, put uniform probabilities on the
candidates in \Cref{lem:bad-node-children}. For large $D$ they satisfy
\eqref{eq:transition-bad}: indeed, $p_y\le c_H^{-1}a^{-D}$ and
$t_y\ge1/8$, while $\gamma,b_*<1$.

\subsection{The capacity estimate}

It remains to construct transitions satisfying
\eqref{eq:transition-general}. A single subdivision depth $D$ must
work even when a node's interval $J$ is arbitrarily short. We therefore
integrate over $J$ before choosing $D$; its length will cancel from
the estimate.

\begin{lemma}[Separated capacity]
\label{lem:separated-capacity}
For sufficiently large $D$, uniformly over all nodes, there is an
$Hr$-separated subset $Y'\subset Y$ such that
\begin{equation}
 \sum_{y\in Y'}\min\{a^{-\gamma D}t_y,a^{\tau D}t_y^2\}\geq1.
 \label{eq:capacity}
\end{equation}
\end{lemma}
\begin{proof}
Join two candidates in a graph when their distance is at most $Hr$.
Packing bounds its degree, so graph coloring partitions $Y$ into at most
$C_H$ separated classes, with $C_H$ independent of $D$. If the conclusion fails, then
\begin{equation}
 \sum_{y\in Y}\min\{a^{-\gamma D}t_y,a^{\tau D}t_y^2\}<C_H.
 \label{eq:failed-capacity}
\end{equation}
Call $y$ \emph{large} when $t_y\geq a^{-(\tau+\gamma)D}$ and
\emph{small} otherwise. On large candidates the minimum is linear, so
\begin{equation}
 \sum_{\text{large }y}t_y\leq C_Ha^{\gamma D}.
 \label{eq:large-total-length}
\end{equation}

We will test levels at an intermediate radius $s$, larger than the
child radius $r$ and smaller than the parent radius $R$. This provides
more radial tests than the large ranges can occupy. At the same time,
it lets us detect the remaining levels with a modulus cost at the
relative scale $r/s$.

\emph{Remove the large ranges.}
Set $D'=\lceil\beta D\rceil$ and $s=a^{-D'}R$. For large $D$,
$0<D'<D$ and $s<R/4$. Choose $s$-spaced radial bins $\mathcal B$
in $[(A+1)R,(B_0-1)R]$, so $\#\mathcal B\asymp a^{D'}$.
On $\mathcal B\times J$ use counting measure times Lebesgue measure.
Mark $(b,t)$ if a large candidate satisfies
\[
 |d(x,y)-b|\leq5s,\qquad t\in I_y^{\mathrm{out}}.
\]
Each candidate meets at most 11 bins. By
\eqref{eq:closed-range-length} and \eqref{eq:large-total-length}, the
marked measure is at most $CC_Ha^{\gamma D}\Delta$. Since $\gamma<\beta$,
the unmarked measure is at least
\begin{equation}
 ca^{D'}\Delta.
 \label{eq:unmarked-mass}
\end{equation}

\emph{Detect the remaining levels by modulus.}
Take a maximal $s$-separated set $\mathcal Z_s$ in $\mathcal A_x(R)$;
its radius-$s$ balls cover that annulus, and it has at most $Ca^{QD'}$ points. Expand at
each $z\in\mathcal Z_s$ at depth $k+D'$ by $h_z$. Put
$q=r/s=a^{-(D-D')}$, and take macroscopic admissible weights
$\omega_\xi$, $\xi\in T_{D-D'}$, of squared mass $M(D-D')$.
For fixed $C_T>2+C_E$, define
\[
 \sigma_y^z=
 \begin{cases}
 \displaystyle\sum_{\xi:\ d(\xi,h_zy)\leq C_Tq}\omega_\xi,
       &d(y,z)\leq3s,\\[3pt]
 0,&d(y,z)>3s,
 \end{cases}
 \qquad \sigma_y=\sum_{z\in\mathcal Z_s}\sigma_y^z.
\]
Within $\overline B(z,3s)$ the candidate images are separated by a
fixed multiple of $q$. Thus every sum has boundedly many terms and
every target weight is charged boundedly many times. Thus
$\sum_y(\sigma_y^z)^2\leq CM(D-D')$ for each $z$. Also each $y$
lies within $3s$ of boundedly many auxiliary centers. Packing and
Cauchy--Schwarz give
\begin{equation}
 \sum_y\sigma_y^2\leq Ca^{QD'}M(D-D').
 \label{eq:capacity-transfer-energy}
\end{equation}

Fix an unmarked $(b,t)$. Choose $w\in C_t$ on the radius-$b$ sphere
and $z\in\mathcal Z_s$ with $d(w,z)\leq s$. The diameter bound gives
a point of $C_t$ outside $\overline B(w,2s)$. Join these points in the
$r/4$-neighborhood of $C_t$, by \Cref{lem:paths}, and stop at the first
radius-$s$ hit about $w$. The resulting path $\pi$ has diameter at
least $s$ and lies in $\overline B(w,s)\subset\overline B(z,2s)$.
The scale-$s$ expansion $h_z$ therefore makes it macroscopic by the
choice of $d_0$ in \Cref{sec:geometry}.

The radial bin lies between $(A+1)R$ and $(B_0-1)R$, while $s<R/4$.
Thus every point of $\pi$ lies between radii $(A+3/4)R$ and
$(B_0-3/4)R$ about $x$. In particular it lies in the candidate
annulus and has a candidate center within distance $r$. The same
fixed source margin places every level point within $r/4$ of $\pi$
in $K^+\subset U$. On these points the scalar-limit property gives
$u=t$, even though it does not assert this identity on parts of
$C_t$ outside $U$.

For a target cell meeting $h_z\pi$, choose a preimage hit $v\in\pi$
and $y\in Y$ with $d(y,v)\leq r$. Then $d(y,z)<3s$. If $\xi$ is the target cell's center,
\eqref{eq:expansion} gives $d(\xi,h_zy)<(2+C_E)q<C_Tq$, so its
weight occurs in $\sigma_y^z$. A point of $C_t$ is within $r/4$ of
$v$, hence within $5r/4<c_0r$ of $y$. The preceding interior check
gives $t\in I_y^{\mathrm{out}}$. Since
$|d(x,y)-b|\leq s+r<5s$, unmarkedness makes $y$ small. Admissibility
therefore yields
\begin{equation}
 \sum_{\substack{\text{small }y:\ |d(x,y)-b|\leq5s}}
 \sigma_y\mathbf1_{I_y^{\mathrm{out}}}(t)\geq1
 \quad\text{on every unmarked pair}.
 \label{eq:unmarked-indicator}
\end{equation}

The unmarked levels now have an admissible detection weight supported
on small candidates. We compare its squared mass with the surviving
amount of level parameter.

\emph{Compare the energies.}
Integrating \eqref{eq:unmarked-indicator}, using bounded bin multiplicity,
gives
\[
 ca^{D'}\Delta\leq C\Delta\sum_{\text{small }y}\sigma_y t_y.
\]
Only finitely many closed-interval indicators are integrated; no
measurable selection of continua or paths is needed, and singleton
ranges contribute zero. Cancel $\Delta>0$. Cauchy--Schwarz and
\eqref{eq:capacity-transfer-energy} imply
\begin{equation}
 \sum_{\text{small }y}t_y^2
 \geq\frac{ca^{-(Q-2)D'}}{M(D-D')}.
 \label{eq:small-square-capacity}
\end{equation}
Subexponential growth gives $M(D-D')\leq a^{\tau D/4}$ for large $D$: first
bound $M(j)\leq Ca^{\tau j/8}$ and then absorb $C$ by increasing $D$. Since $D'\leq\beta D+1$ and $(Q-2)\beta<\tau/4$,
\[
 \sum_{\text{small }y}a^{\tau D}t_y^2
 \geq ca^{[3\tau/4-(Q-2)\beta]D}
 \geq ca^{\tau D/2}.
\]
On small candidates this is exactly the capacity, contradicting
\eqref{eq:failed-capacity}. Every threshold for $D$ is uniform over nodes.
\end{proof}

Fix such a $D$. At each good node, divide the positive capacities on a
set $Y'$ from \Cref{lem:separated-capacity} by their sum. These
probabilities satisfy \eqref{eq:transition-general}. At the remaining
nodes use the uniform rule above. Every positive transition retains
a nonempty open child interval $J_y\subset J$.

\subsection{Common target patterns}

The probabilities are now fixed. We verify the additional geometric
information carried by a good node: its two source variations expand
into a pair of balls chosen from a finite list, with separated range
intervals. The list depends on the fixed $D$, but not on the node's
generation or interval length.

Fix a node's expansion $h=v^{-1}$ when constructing it. At a good node,
join the extrema of each inner ball by a path in
$B(y^\pm,c_{\mathrm{in}}r)$. Each path has endpoint variation at least
$\kappa\Delta$. Its points lie at distance at least $(A-2)R$ from $x$,
so \eqref{eq:expansion} puts its image within
$c_{\mathrm{pat}}\theta/2$ of $hy^\pm$.
Conversely, write $D_w=\max\{R,d(x,w)\}$. If $d(y,z)\geq c_0r$, then
$D_z\leq D_y+d(y,z)$ and monotonicity of $s/(D_y+s)$ give
\[
 d(hy,hz)\geq\frac{c_ERc_0r}{D_y(D_y+c_0r)}
 \geq\frac{c_Ec_0\theta}{(B_0+2)^2}>12c_{\mathrm{pat}}\theta.
\]
Thus
\begin{equation}
 \overline B(hy,10c_{\mathrm{pat}}\theta)
       \subset h(B(y,c_0r)).
 \label{eq:pattern-containment}
\end{equation}
Choose a fixed finite $c_{\mathrm{pat}}\theta$-net in $Z$. For each
$hy^\pm$ choose a net point within $c_{\mathrm{pat}}\theta$ and take
the closed radius-$3c_{\mathrm{pat}}\theta$ ball about it. The ordered
pair $(\mathcal B^-,\mathcal B^+)$ is the node's \emph{pattern}.
Each ball contains its path image and lies in the corresponding outer
source-ball image. Consequently, for $u\circ h^{-1}$ the two whole-ball
range intervals have distance at least $\Delta/4$, and each ball
contains a path with endpoint variation at least $\kappa\Delta$.
There are finitely many patterns, independently of depth.

\subsection{The tree and its flow}

Choose $k_0$ so large that $R_0=r_{k_0}$ satisfies all the source-scale
conditions above and
\[
 2(B_0+1)R_0<\epsilon_0/2,\qquad (B_0+2)R_0<d_1/2.
\]
Non-local-constancy at $x_0$ supplies a nonempty open interval
$J_0\subset(0,1)$ strictly inside the range on
$\overline B(x_0,c_0R_0)$. Write $\Delta_0=|J_0|$ and construct the
tree recursively, storing each node's expansion and, at good nodes,
its pattern.

A child moves by at most $(B_0+1)R$ from a parent of scale $R$.
Consequently all descendants stay within
\begin{equation}
 \frac{B_0+1}{1-\theta}R\leq2(B_0+1)R
 \label{eq:descendant-displacement}
\end{equation}
of that parent, and all centers lie in $K$. Generation $n$ has scale
$r_{k_0+Dn}$ and its centers are $H$ times that scale apart: this holds
for siblings, while children of distinct parents of scale $R$ are
separated by at least
\[
 [H-2(B_0+1)]R>HR/2>H\theta R.
\]
At any level of scale $R$, \eqref{eq:measure} bounds $\Delta_i^2$ by
$C\mu(\overline B(x_i,C_\mu c_0R))$. These enlarged balls are disjoint
because $H>4C_\mu c_0$, and hence
\begin{equation}
 \sum_i\Delta_i^2\leq E_0,\qquad E_0=C\mu(Z)<\infty.
 \label{eq:level-energy}
\end{equation}

Give the root probability 1 and let $P_i$ be the product of transitions
to node $i$. Every level has total probability 1. If $i$ has generation
$n$, let $m_i$ count its ancestors at generations $0,\ldots,n-1$ that
are not good. Multiplying the transition bounds and telescoping the
interval ratios yields
\begin{equation}
 P_i\leq a^{-\gamma Dn}\frac{\Delta_i}{\Delta_0},\qquad
 \Delta_i^2\geq\Delta_0^2P_i
 a^{D((b_*+\tau)m_i-\tau n)}.
 \label{eq:flow}
\end{equation}
Conditionally, if $i$ is $m$ generations below $j$, nesting gives
\begin{equation}
 \frac{P_i}{P_j}\leq a^{-\gamma Dm}\frac{\Delta_i}{\Delta_j}
 \leq a^{-\gamma Dm}.
 \label{eq:conditional-flow}
\end{equation}

\begin{lemma}[Levels with substantial good mass]
\label{lem:good-levels}
Infinitely many generations have total probability at least $1/2$
on good nodes.
\end{lemma}
\begin{proof}
By \eqref{eq:flow}, \eqref{eq:level-energy}, and Jensen's inequality,
\[
 \frac{E_0}{\Delta_0^2}
 \geq\sum_iP_i a^{D((b_*+\tau)m_i-\tau n)}
 \geq a^{D((b_*+\tau)\sum_iP_im_i-\tau n)}.
\]
Thus
\[
 \frac1n\sum_iP_im_i\leq\frac{\tau}{b_*+\tau}+O(1/n)
 =\frac1{11}+O(1/n).
\]
Conservation of probability identifies the left side with the average
bad mass over the first $n$ generations. If good mass were eventually
less than $1/2$, that average would have lower limit at least $1/2$,
a contradiction.
\end{proof}

\begin{proposition}[Properties of the flow]
\label{prop:probability-flow}
The construction gives a fixed $D$ and a tree with generation-$n$
centers $Hr_{k_0+Dn}$-separated and descendant displacement
\eqref{eq:descendant-displacement}. Its nonempty value intervals are
nested, its probabilities satisfy \eqref{eq:flow} and
\eqref{eq:conditional-flow}, and infinitely many levels have good mass
at least $1/2$. Every node has a fixed expansion. Each good node has
a pattern from a finite list of ordered target-ball pairs. For that expansion,
each target ball contains a path with endpoint variation at least
$\kappa\Delta_i$, their whole-ball range intervals have distance at
least $\Delta_i/4$, and the balls lie in the images of the respective
outer balls of radius $c_0\theta r_{k_0+Dn}$.
\end{proposition}
